% Section content template for individual Educative lessons
% This file contains the actual content for a specific section
% Generated from Educative API section content

% Section: Code of Jigsaw Puzzle
% Section ID: 4527257987317760
% Section Slug: code-of-jigsaw-puzzle
% Generated: 2025-12-14T13:10:56.129239

We've reviewed the different aspects of the jigsaw puzzle and observed the attributes attached to the problem using various UML diagrams. Now , let's explore the more practical side of things, where we will work on implementing the jigsaw puzzle using multiple languages. This is usually the last step in an object-oriented design interview process.
\\
We have chosen the following languages to write the skeleton code of the different classes present in the jigsaw puzzle:

\begin{itemize}
\item
\protect\phantomsection\label{mFJcZrlVIAIYPcgTd5uW7}
Java
\item
\protect\phantomsection\label{51963932twQ2rjqOuumFr}
C\#
\item
\protect\phantomsection\label{GAgCNG4XUgP-Tf4CK-qwV}
Python
\item
\protect\phantomsection\label{7Usalcza3xyWDLobHHCSJ}
C++
\item
\protect\phantomsection\label{ABeGXGdtC3tNeXK6Zl4S7}
JavaScript
\end{itemize}

\subsection{Jigsaw puzzle}\label{nySktjHQfz4_kAZzIOt8G}
\\
This section will provide the skeleton code of the classes designed in the class diagram lesson.
\begin{quote}
\textbf{Note:} For simplicity, we are not defining getter and setter functions. The reader can assume that all class attributes are private, accessed through their respective public getter methods, and modified only through their public method functions.
\end{quote}
\subsubsection{Enumerations}\label{FK3Cg7Ck_xXgfE9n5k-cQ}
\\
The following code defines the \texttt{Edge} enum that represents the shape of the puzzle piece:
\begin{quote}
\textbf{Note:} JavaScript does not support enumerations, so we will use the \texttt{Object.freeze()} method as an alternative that freezes an object and prevents further modifications.
\end{quote}

\textbf{Enum definitions}

\begin{lstlisting}[language=Java]
enum Edge {
  INDENTATION, 
  EXTRUSION, 
  FLAT
}
\end{lstlisting}

\subsubsection{Side}\label{_ahFF7YPkyYYRp-o1OVUm}
\\
The \texttt{Side} class contains a reference to the \texttt{Edge} enum that will describe the edge shape of any side of a puzzle piece. Its definition is given below:

\textbf{The Side class}

\begin{lstlisting}[language=Java]
public class Side {
  private Edge edge;
}
\end{lstlisting}

\subsubsection{Piece}\label{B6SjSJlx9bFwLHg9WYeif}
\\
The \texttt{Piece} class contains a list of sides (representing each side of a piece) and ensures that a particular piece is a corner, middle, or edge using its respective functions. The definition of this class is provided below:

\textbf{The Piece class}

\begin{lstlisting}[language=Java]
public class Piece {
  private List<Side> sides = Arrays.asList(new Side[4]);
  public boolean checkCorner() {}
  public boolean checkEdge() {}
  public boolean checkMiddle() {}
}
\end{lstlisting}

\subsubsection{Puzzle}\label{ePhNw9GWSys3S70Fvesd0}
\\
The \texttt{Puzzle} class includes the entire board of the jigsaw puzzle and the unused pieces. As the board is a rectangle, it will be represented using a 2D array. In addition, there will be only one instance of the puzzle board in the jigsaw puzzle. Therefore, the \texttt{Puzzle} class will be a Singleton class to ensure that only one instance of the board is created in the entire system. The definition of this class is given below:

\textbf{The Puzzle class}

\begin{lstlisting}[language=Java]
public class Puzzle {
  private List<List<Piece>> board;
  private List<Piece> free; // represents the currently free pieces (yet to be inserted in board)
  public void insertPiece(Piece piece, int row, int column) {};

  // Puzzle is a singleton class that ensures it will have only one active instance at a time
  private static Puzzle puzzle = null;
    
  // Created a static method to access the singleton instance of Puzzle class
  public static Puzzle getInstance() {
      if (puzzle == null) {
          puzzle = new Puzzle();
      }
      return puzzle;
  }  
}
\end{lstlisting}

\subsubsection{Puzzle solver}\label{H8I77m23rmlkOaAizdmDE}
\\
The \texttt{PuzzleSolver} class tries to solve the puzzle using the \texttt{matchPieces()} function. Its definition is given below:

\textbf{The PuzzleSolver class}

\begin{lstlisting}[language=Java]
public class PuzzleSolver {
  public Puzzle matchPieces(Puzzle board) {}
}
\end{lstlisting}

\subsection{Wrapping up}\label{oduvT-U5GDbYO8Cu-X2TP}
\\
In this chapter, we've explored the complete design of a jigsaw puzzle. We 've looked at how a basic jigsaw puzzle can be visualized using various UML diagrams and designed it using object-oriented principles and design patterns.

% End of section content